\summarypart{SEVEN}{The Age of Progress}{1789--1914}
\begin{partoverview}
To the intellectual challenges that came with modern science, this age added social unrest. The French
Revolution released new hopes and expectations for the common man, and power now seemed within reach
of the masses. How were Christians to meet the needs of the crowded cities? Was man simply a product of
evolutionary forces? Christians were seriously divided over the answers. Without the traditional
support of the state, many Protestants turned to voluntary societies, both to minister to the poor and
the oppressed and to carry the gospel to foreign lands.
\end{partoverview}
\partchapters{%
  \pch{36}{The Restoration of Fortresses}{Catholicism in the Age of Progress}
  \pch{37}{A New Social Frontier}{Nineteenth-Century England}
  \pch{38}{To Earth's Remotest People}{Protestant Missions}
  \pch{39}{The Destiny of a Nation}{A Christian America}
  \pch{40}{A Bridge for Intelligent Moderns}{Protestant Liberalism}
  \pch{41}{Nothing to Lose but Chains}{The Social Crisis}}
